\documentclass[11pt,a4paper]{article}

% ── Packages (only those available on this installation) ──────────────────────
\usepackage[margin=2.5cm]{geometry}
\usepackage{amsmath,amssymb}
\usepackage{graphicx}
\usepackage{booktabs}
\usepackage{longtable}
\usepackage{array}
\usepackage{xcolor}
\usepackage{hyperref}
\usepackage{caption}
\usepackage{subcaption}
\usepackage{enumitem}
\usepackage{microtype}
\usepackage{listings}
\usepackage{float}
\usepackage{placeins}

% ── Figure path ───────────────────────────────────────────────────────────────
\graphicspath{{figures/}}

% ── Hyperref setup ────────────────────────────────────────────────────────────
\hypersetup{
  colorlinks = true,
  linkcolor  = blue!70!black,
  citecolor  = green!50!black,
  urlcolor   = blue!70!black,
}

% ── Custom commands ───────────────────────────────────────────────────────────
\newcommand{\numu}{$\nu_\mu$}
\newcommand{\numubar}{$\bar{\nu}_\mu$}
\newcommand{\ccpip}{$\mathrm{CC}\pi^{\pm}$}
\newcommand{\pmu}{$p_\mu$}
\newcommand{\ppi}{$p_\pi$}
\newcommand{\costhmu}{$\cos\theta_\mu$}
\newcommand{\costhpi}{$\cos\theta_\pi$}
\newcommand{\thmupi}{$\theta_{\mu\pi}$}
\newcommand{\uboone}{MicroBooNE}
\newcommand{\GeVc}{\,\text{GeV}\!/c}
\newcommand{\pot}{\,\text{POT}}
\newcommand{\cm}{\,\text{cm}}

% inline math shorthand so they work inside math mode too
\newcommand{\pmuM}{p_\mu}
\newcommand{\ppiM}{p_\pi}
\newcommand{\cthmuM}{\cos\theta_\mu}
\newcommand{\cthpiM}{\cos\theta_\pi}
\newcommand{\thmupM}{\theta_{\mu\pi}}

% ── Listings (C++ snippets) ───────────────────────────────────────────────────
\lstset{
  language=C++,
  basicstyle=\ttfamily\small,
  keywordstyle=\bfseries\color{blue!70!black},
  commentstyle=\itshape\color{gray},
  backgroundcolor=\color{gray!10},
  frame=single,
  framerule=0.4pt,
  numbers=none,
  breaklines=true,
}

% ── Table column helpers ──────────────────────────────────────────────────────
\newcolumntype{C}[1]{>{\centering\arraybackslash}p{#1}}
\newcolumntype{R}[1]{>{\raggedleft\arraybackslash}p{#1}}
\newcolumntype{L}[1]{>{\raggedright\arraybackslash}p{#1}}

% Let TeX stretch spacing on hard-to-break lines (long \texttt paths/code) to keep
% text out of the margin.
\emergencystretch=3em

% ─────────────────────────────────────────────────────────────────────────────


\begin{document}

\begin{titlepage}
\centering
\vspace*{2cm}
{\LARGE\bfseries Change Log\par}
\vspace{0.5cm}
{\large Measurement of $\nu_\mu + \bar{\nu}_\mu$ Charged-Current Single Pion
Production on Argon in the NuMI Beam\par}
\vspace{0.8cm}
{\large Corrections that changed a result\par}
\vspace{0.5cm}
{\normalsize Draft 1.4, 2026-09-04\par}
\vspace{0.15cm}
{\footnotesize\ttfamily release 2026-09-04 \textperiodcentered\ 51 differential extractions $+$ 6 one-bin totals \textperiodcentered\ blind\par}
\vspace{1.5cm}
{\large Ishanee Pophale, Jarek Nowak\par}
\vspace{0.5cm}
{\normalsize\textit{Lancaster University}\par}
\vspace{2cm}

\begin{abstract}
This document records every correction that materially moved a number in the
$\mathrm{CC}\pi^{\pm}$ NuMI analysis, with the value before and after.  It exists so
that a reader of the analysis note never has to work out whether a given table predates
a given fix, and so that a result quoted from an older draft can be identified as
superseded rather than merely different.

Most entries are normalisation errors.  That concentration is itself worth stating: the
extraction divides by four numbers --- the data exposure, the integrated flux, the argon
target count, and the per-run Monte Carlo exposure --- and eight separate defects came
from one of those four differing between configurations without anything detecting it.
The response is \texttt{norm\_manifest.py}, which reads the authoritative configuration
files and asserts that they agree, so that a future drift fails a check instead of
producing a wrong cross section.

Two corrections changed physics conclusions rather than only normalisation, and the
conclusions they invalidated are marked as withdrawn here and in the analysis note: the
non-uniform-grid regularisation defect, which had collapsed several additional-smearing
matrices towards rank one, and the beam-axis convention, under which reconstructed
angular and transverse observables were not the same quantities as the generator
predictions they were being compared against.
\end{abstract}
\end{titlepage}

\tableofcontents
\clearpage

\section{How to read this document}

Each entry states what was wrong, how it was found, what changed numerically, and what
--- if anything --- had to be withdrawn.  ``Before'' values are quoted from the draft in
which they appeared and should not be used; they are given so that an older table can be
recognised.

Entries are ordered by the stage of the analysis they affect, not chronologically:
exposure and flux first, then the Monte Carlo weighting, then the unfolding, then the
systematic covariance.

\section{Exposure, flux, and target normalisation}

\subsection{Obsolete integrated flux (factor ${\sim}3.48$)}

An outdated flux file was used to convert the extracted event rate into a cross section.
The integrated flux it supplied was larger than the correct value by a factor of
approximately $3.48$, so every cross section produced with it was smaller than the truth
by the same factor.

\textbf{Effect.} All cross sections, all configurations.  Any absolute normalisation from
a draft predating this fix is wrong by ${\sim}3.48$ and must not be quoted.  Shapes were
unaffected, because the error was a single multiplicative constant.

\subsection{Configuration-specific flux was not being used}

The RHC and combined extractions were normalised with the FHC integrated flux.  The three
exposures see genuinely different spectra, and the correct integrated fluxes are
\[
\Phi_{\mathrm{FHC}} = 6.81159\times10^{-10},\quad
\Phi_{\mathrm{RHC}} = 6.44646\times10^{-10},\quad
\Phi_{\mathrm{comb}} = 6.60865\times10^{-10}
\]
in $\mathrm{cm}^{-2}\,\mathrm{POT}^{-1}$.

\textbf{Effect.} RHC cross sections rose by $5.7\%$ and combined by $3.1\%$.  FHC was
already correct and did not move.  The flux is now set by an explicit \texttt{Flux} line
in each cross-section configuration file rather than inherited, and
\texttt{norm\_manifest.py} checks that all configurations of a given exposure agree on it.

\subsection{Per-mode rather than per-run exposure weighting}

Data and Monte Carlo exposures were pooled per horn mode and a single scale factor
applied.  Because the data-to-Monte-Carlo exposure ratio differs substantially between
runs, this weighted the result by how much simulation happened to exist for each run
rather than by the data actually collected.  Each run is now scaled by its own
$S_r = \mathrm{POT}^{\mathrm{data}}_r/\mathrm{POT}^{\mathrm{MC}}_r$.

\textbf{Effect.} Changed the relative weight of runs within every configuration.  This is
also the change that makes the framework ready for real data, where the per-run exposures
are not free parameters.

\subsection{Duplicated run-total POT in split files}

Two conventions coexist in the processed samples.  In \texttt{processed/} a run split
across several files repeats the \emph{run-total} exposure in every file; in
\texttt{processed/w/} each file carries its own \emph{true} exposure.  Dividing by the
per-file value broke the first convention and summing over files broke the second: the
first made the inclusive tune come out a factor of two low, the second drove an RHC
proton-tagged cross section negative.

The resolution is to normalise by the sum of \emph{distinct} per-run exposures, which is
correct under both conventions.

\textbf{Effect.} At the time of the fix the RHC $W_{\pi p}$ cross section moved from
$-0.194$ to $+0.578$ (in units of $10^{-38}\,\mathrm{cm}^2/\mathrm{Ar}$), and the
inclusive RHC $p_\mu$ result was unchanged at $0.860$.  Results retracted on the strength
of the negative value were reinstated.  Both numbers have since moved again under the
central-value weighting fix below; the current values are in the analysis note.

\section{Monte Carlo weighting}

\subsection{Central-value weight applied twice to the fake data}

This is the largest single correction and it affected every result.

The fake data are generated externally by \texttt{throw\_perrun\_*.C}, which folds the
tuned central value, the PPFX central value and the normalisation into the event
multiplicity and then writes the per-event weights back as unity --- the throw is a
Poisson realisation of an already-weighted expectation.  The framework's
\texttt{build\_universes} nevertheless substituted the \emph{weighted} central-value
histogram for the raw counts, on both the reconstructed and the truth side, so the
central value was applied a second time.

\textbf{How it was found.} By asking directly whether the unfolded result was being
compared against the input that generated it.  Four independent checks confirmed the
diagnosis: $1621.0$ is a direct count of selected events across the four FHC files;
$1721.9$ equals both \texttt{FakeDataMC\_0\_reco} and
$\texttt{onBNB\_reco}-\texttt{extBNB\_reco}$; run~1 gives $558.0$ entries unweighted
against $599.9$ weighted for the same $558$ events; and the throw's own expectation of
$1567.1$ agrees with the thrown count at $1.3\sigma$.

\textbf{Effect.} The mean inclusive closure ratio moved from $1.106$ to $1.001$ ($1.016$ once the truth-side correction below is included, $1.018$ on the re-thrown fake data), and all
$51$ extracted cross sections fell by between $3.8\%$ and $12.9\%$.  Every absolute cross
section from a draft predating this fix is high by roughly a tenth.

\textbf{Withdrawn.} Two earlier explanations of the closure excess --- one attributing it
to the unfolding covariance and Wiener filter, one to $A_C$ smearing --- were both wrong
and are retracted.  They are retained in the technical supplement as recorded dead ends
rather than deleted.

\subsection{Generator prediction histograms truncated at the overflow}

Generator predictions were read without their overflow contents, so the highest-momentum
region was silently missing from the comparison curves.

\textbf{Effect.} Model comparison only; the measured cross sections were unaffected.  Any
generator $\chi^2$ from a draft predating the regeneration is invalid.

\section{Unfolding}

\subsection{Second-derivative regularisation on a non-uniform grid}

The Wiener-SVD second-derivative regulariser used the uniform-grid $(1,-2,1)$ stencil on
bins whose widths differ by factors of ten to fifteen.  On such a grid that operator is
not an approximation to the second derivative of anything, and its near-null direction is
not the physically flat one.  The Wiener filter can then retain that direction alone,
which drives the additional-smearing matrix $A_C$ towards rank one and destroys the shape
information the measurement is supposed to carry.

The corrected operator differences the \emph{density} on a non-uniform grid, reducing to
the old stencil when the widths are equal, with one-sided first-derivative rows at the
ends so that the operator stays non-singular without asserting that the distribution
vanishes at the edges --- which it does not, since $\cos\theta_\mu$ peaks at the upper
edge.

\textbf{Effect} on the conditioning measure $s_2/s_1$ of $A_C$:
\begin{center}
\begin{tabular}{lcc}
\toprule
Observable & before & after \\
\midrule
FHC $\theta_\mu$              & $0.015$ & $0.51$ \\
combined $\theta_\mu$         & $0.045$ & $0.38$ \\
proton-tagged $W_{\pi p}$     & $0.022$ & $0.60$ \\
proton-tagged $p_n$           & $0.21$  & $0.88$ \\
\bottomrule
\end{tabular}
\end{center}

\textbf{Withdrawn.} The earlier conclusion that $\theta_\mu$ ``repairs'' what
$\cos\theta_\mu$ cannot is retracted: under the defective operator every model was being
collapsed onto a common shape, so the apparent agreement was an artefact of the
regularisation rather than a property of the data.  After the fix $\theta_\mu$
discriminates between models again and prefers the tune that generated the fake data.

This entry is the clearest illustration of why closure alone is not a sufficient
validation.  A rank-one extraction closes perfectly against its own correspondingly
smeared truth while retaining almost no shape information.

\subsection{Angular and transverse observables were defined about detector \texorpdfstring{$z$}{z}}

The NuMI neutrinos arrive approximately $28.6^\circ$ away from the detector $z$ axis, but
reconstructed angles and transverse kinematic imbalance were being computed about $z$
while the generator predictions were computed about the neutrino direction.  These are
not the same observables.  Truth-level quantities now use the per-event neutrino
direction and reconstruction uses the target-to-vertex direction; the two constructions
differ only at the milliradian level.

\textbf{Effect.} $64.6\%$ of signal events change $\cos\theta_\mu$ bin and $44.6\%$ change
$\cos\theta_\pi$ bin.  The mean $\delta p_T$ moves from $0.742$ to
$0.264\ \mathrm{GeV}/c$.  The muon and pion momenta, the opening angle $\theta_{\mu\pi}$,
the hadronic mass $W_{\pi p}$, and the event selection are all unaffected.

\textbf{Withdrawn.} All proton-tagged transverse-kinematic model comparisons made before
this correction are withdrawn, because the measured and predicted quantities were
different observables.

\subsection{Predictions were smeared by \texorpdfstring{$A_C$}{A C} twice}

Every model prediction plotted or compared against the data was multiplied by the
additional-smearing matrix twice: once by the extractor, which applies $A_C$ to each entry
of the prediction map in place, and again on the plotting path.  Confirmed exactly, as
$A_C^{2}\cdot\text{raw} = \text{plotted}\times\text{width}$ to $1.0000$ in every bin.

\textbf{How it was found.}  By testing whether the released $A_C$ lets an external user
reproduce a published generator curve.  It did not --- the discrepancy was $31\%$ on the
integral and varied from $1.21$ to $1.55$ across bins, so it was not a normalisation --- and
the residual factor turned out to be $A_C$ itself.

\textbf{Effect.}  Model $\chi^2$ values only.  The anomaly the note previously attributed
to $A_C$ suppression of sharply peaked generators was this bug:

\begin{center}
\small
\begin{tabular}{llcc}
\toprule
Observable & Model & before & after \\
\midrule
FHC $\theta_\mu$      & NEUT  & $40.60/5$ & $8.96/5$ \\
FHC $\theta_\mu$      & NuWro & $38.62/5$ & $9.43/5$ \\
FHC $\cos\theta_\mu$  & NEUT  & $29.78/5$ & $4.90/5$ \\
FHC $\cos\theta_\mu$  & NuWro & $29.10/5$ & $4.93/5$ \\
FHC $p_\mu$           & GENIE & $5.52/7$  & $0.57/7$ \\
\bottomrule
\end{tabular}
\end{center}

\textbf{Not affected.}  The unfolded cross sections (FHC $p_\mu$ integrates to $0.8275$
before and after) and the closure test, because the fake-data truth was the one curve not
passed through the second multiplication.  The previous table also predated the
central-value weighting fix, so part of the shift in its ``truth'' column is that
correction rather than this one.

\textbf{Withdrawn.}  The argument for preferring $\theta$ over $\cos\theta$ on the grounds
that the $\cos\theta_\mu$ $\chi^2$ was an artefact.  The effect that motivated it does not
exist.

\section{Systematic covariance}

\subsection{Combined detector-variation exposure double-counted}

The combined configuration pooled detector-variation samples per horn mode, and the
denominator used a single file's exposure rather than the summed Monte Carlo exposure over
all files of that variation type, so the combined detector systematic was inflated.

\textbf{Effect.} The total detector systematic fell from $42.3\%$ to $24.6\%$, and the
combined result moved from $0.910$ to $1.027$ times its expected value.  FHC and RHC were
bit-identical before and after, which is the check that isolates the defect to the
combination step.

\subsection{Proton-tagged detector systematics were absent entirely}

Of the fifty-one extractions, the thirty-three proton-tagged ones carried \emph{no}
detector systematic at all: the \texttt{detVar} column was \texttt{n/a} for every
proton-tagged observable in every configuration, so their quoted uncertainties omitted a
term worth $17$--$43\%$ in the corresponding inclusive results.

The cause was a configuration gap rather than missing data.  The proton-tagged detector
variation samples had been processed on 2026-08-20/21 and carry the required branches, but
nothing referenced them: the \texttt{\_w} file-properties tables contained no
\texttt{detVar} lines, and \texttt{ccpi1p\_systcalc\_numi.conf} had no \texttt{DV} block
at all.  Both were needed; adding either alone does nothing.

\textbf{Effect.}  The thirty-three proton-tagged universes were rebuilt with the nine
detector variation samples included and re-unfolded from them.  Every extraction in the
analysis now carries a detector term.  The proton-tagged terms run $10.5\%$--$64.4\%$
with a median of $23.9\%$, against $16.9\%$--$42.8\%$ for the inclusive results.

\textbf{Check.}  Per-mode detector-variation double counting has occurred in this analysis
before, and the combined configuration --- which carries two per-mode sets where the
single-mode configurations carry one --- is the one exposed to it.  Double counting
inflates the combined result.  Of the eight observables (of eleven) whose combined value
falls outside its own FHC--RHC range, six fall \emph{below} both and only two above, which
is the finite-sample averaging-down expected of statistically limited detector-variation
samples rather than a recurrence.

\subsection{Cut-and-count uncertainty quoted from a pre-rebuild breakdown}

The total flux-averaged cross sections of the cut-and-count path took their systematic
fraction from the $p_\mu$ systematic-breakdown tables, which had not been regenerated
after the detector-variation rebuild of 2026-08-31 (their detector terms were
$\sim10\%$ against $17$--$43\%$ in the release).  The breakdown tables in the note and the
supplement were regenerated from the release systematics dumps on 2026-09-02 and the
cut-and-count uncertainties recomputed from them; the macro also now applies the
central-value MC weight, which moved the central values by $0.3$--$0.5\%$.

\textbf{Effect} (values of the time; the current totals are in the one-bin entry below).  Inclusive totals ($10^{-38}$~cm$^2$/Ar): FHC $1.06\pm0.27\to1.03\pm0.43$,
RHC $0.98\pm0.35\to0.93\pm0.42$, combined $1.02\pm0.30\to0.98\pm0.40$.  The generator
comparison changes from ``GENIE $\sim1\sigma$ low'' to ``GENIE $0.5$--$0.6\sigma$ low''.
A proton-tagged cut-and-count total is quoted for the first time (FHC $0.63\pm0.31$,
RHC $0.55\pm0.21$, combined $0.59\pm0.25$).

\subsection{Fake-data truth still carried the central-value weight}

The 2026-08-30 correction (above) stopped the tune weight being applied a second time
to the externally thrown fake data on the reconstructed side.  The truth side was
missed: the throw macro resets only the scalar weight branches, the
\texttt{weight\_TunedCentralValue\_UBGenie} vector survives, so the universe maker still
writes a tune-weighted ``CV'' universe for each fake-data file, and the systematics
calculator took that weighted histogram as the fake-data \emph{truth} whenever it was
present.  Measured on the FHC $p_\mu$ universe file: weighted/unweighted true-signal
integrals $1.023$, $1.008$, $1.008$, $1.018$ for Runs~1, 2, 4, 5, about $1.6\%$ overall.
Every closure ratio was therefore low by that amount and every closure $\chi^2$ slightly
too small.  Found in the 2026-09-02 code trace; the truth path now uses the same
external-throw flag as the reco path, and the universe cache key was bumped so the
stale sums cannot be reused.

\textbf{Effect.}  Unfolded cross sections and covariances are unchanged (FHC $p_\mu$
integrates to $0.8275$ before and after).  Closure: FHC $p_\mu$ unfolded/truth
$1.024\to1.041$, $\chi^2$ $0.10/7\to0.13/7$; the mean inclusive closure ratio
$1.001\to1.016$.  All fifty-one sidecars, the released curves
(\texttt{truth\_smeared} column), the closure tables and the closure figures were
regenerated; the ensemble bias is unaffected because its reference is the fixed
central-value truth, not the thrown one.

\subsection{Beam-off gate count was an event count}

Every live file-properties configuration scaled the pooled Run-1$+$Run-3b beam-off
sample by $\text{bnb\_trigs}[\text{run}]/3\,821\,593$.  The denominator was inherited from
the earlier custom selection, where it was the sum $610\,496+3\,211\,097$ of the two
samples' \emph{event} counts; their gate counts, from the analyser's trigger table
supplied on 2026-09-02, are $4\,582\,248.27+32\,649\,128.65=37\,231\,376.92$, a factor
$9.74$ larger.  The selected cosmic rate per gate exposed it: $5.8$ per $10^{6}$ gates for
the pooled sample against $0.37$--$1.3$ for the dedicated per-period files, versus $0.59$
with the gate count.  Corrected in all six live configurations and in the counting and
cut-flow macros; the universes were re-summed (the cache key tracks the file-properties
contents) and all fifty-one extractions, covariances and the release were regenerated.

\textbf{Effect.}  EXT background at the FHC final cut $128\to13$ events (RHC
$133\to14$, combined $261\to27$); purity $55.8\to59.7\%$ (FHC).  Fake-data central values
are unchanged by construction (the same EXT is added and subtracted); the EXT-statistics
and data-statistics covariance terms shrink, which moves the Wiener-SVD results only
through the regularisation (see the regenerated tables).  A real-data extraction with the
old value would have over-subtracted $\approx\!115$ events, $\approx\!12\%$ of the FHC
signal.  \textbf{Withdrawn:} every beam-off yield, purity and EXT-statistics term quoted
in drafts before 1.2.  The sideband tables were regenerated with the per-run treatment
(\texttt{regen\_sidebands.sh}); the cosmic control region, formerly EXT-dominated by
$4.5:1$, now has $0.7$ EXT events per neutrino-MC event.

\subsection{Beam-off matched by run period instead of one pooled sample}

Beam-off data are cosmic-only, so a sample is matched to a beam-on run by its run period,
irrespective of the horn polarity it was recorded under (analyser, 2026-09-02).  The
pooled Run-1$+$Run-3b sample was replaced by the dedicated per-period samples (Run~1;
Run~3b; Run~4a, 4b, 4c, 4d pooled; Run~5; the Run-1 sample standing in for Run~2), each
scaled by its own gate count from the analyser's table, and processed with the swtrig
fail-open build so that the CRT-era samples select cosmics at all.  The selected rate per
$10^{6}$ gates is $1.75$ (Run~1), $0.43$ (Run~3b), $1.15$ (Run~4), $0.52$ (Run~5), so
period matching moves the prediction by a factor of about two relative to the pooled
sample.  All fifty-one universe files were rebuilt (SLURM array 3205007, 51/51), the
extractions re-unfolded, covariances re-harvested and the release regenerated.

\textbf{Effect.}  EXT background at the final cut (FHC / RHC / combined): pooled
$13/14/27$ $\to$ per-run $30/23/53$ events, i.e.\ $1.9\%$ of the selected sample; purity
$59.1\%$ (FHC).  EXT-statistics term on FHC $p_\mu$ $0.4\to0.9\%$.  Unfolded central
values move by $\lesssim1\%$ through the regularisation.

\subsection{PPFX central-value weight absent from the cut-flow and cut-and-count counts}

The framework's central-value event weight for NuMI is tune $\times$ PPFX~CV $\times$
normalisation (\texttt{UniverseMaker.cxx}), and the per-run fake data are thrown with
that product.  The selection's cut-flow histograms and the cut-and-count macro carried
the tune only.  The PPFX central-value weight averages $0.944$ on selected signal, so
every absolute yield in the cut-flow table and every cut-and-count total sat $3$--$6\%$
above the prediction the unfolding uses; efficiencies and purities, being ratios, were
unaffected at the $0.5\%$ level.  Found in the 2026-09-02 code trace.

\textbf{Effect.}  Cut-and-count totals: FHC $1.06\to1.03$, RHC $0.99\to0.93$, combined
$1.02\to0.98$ ($10^{-38}$~cm$^2$/Ar), and the RHC/FHC ratio $0.928\to0.901$.  The
selection now fills the cut-flow histograms with the full weight and the fourteen
overlay files and the dirt overlay were reprocessed on 2026-09-03
(\texttt{reprocess\_cutflow\_cf.sh}, into a separate directory so the universes' POT
convention is untouched); the cut-flow table is regenerated from them.  No unfolded
result is affected.

\subsection{Proton-tagged TKI figure drawn from the withdrawn fine-binned sidecars}

The proton-tagged W/TKI montage macros read closure sidecars keyed \texttt{dpt},
\texttt{dalphat}, \texttt{dphit}, \texttt{pn}: the withdrawn fine-binned extraction of
2026-08-18, which predates the beam-axis correction, rather than the released two-bin
\texttt{*2bin} files.  The figure therefore compared transverse quantities about detector
$z$ (its ``data'' and tune) with beam-frame generator curves, and showed shape
disagreements of a factor of several that do not exist in the release.  The staleness
check maps figures by configuration only and could not catch it.  Macros repointed, the
stale sidecars quarantined, the figure regenerated (2026-09-03).

\textbf{Effect.}  Figure only; no released number.  In the two-bin scheme the generators
differ by normalisation ($\times1.6$ GENIE to NEUT) with shapes consistent with the data.

\subsection{Result tables rewritten into the wrong place by the table patcher}

The script that repatched the result tables after each regeneration located a table by
the \texttt{\textbackslash begin\{tabular\}} nearest its label, falling back to the
\emph{preceding} tabular whenever the caption between label and tabular exceeded 400
characters.  Four tables were affected in the committed Draft~1.2: the FHC results table
received the RHC rows, the RHC table the combined rows, the combined table the rows of the
integrated-cross-section summary, and in the supplement the cut-and-count comparison was
overwritten by an FHC systematics breakdown.  Caught by the second external review
(2026-09-03).  All release-derived tables are now regenerated as whole tabulars by
\texttt{rebuild\_tables.py} with a locator that requires the tabular to follow the label
inside the same environment, and \texttt{report/check\_tables.py} asserts, on every build,
that each table sits under its own label, has the column count its header declares, and
agrees with \texttt{current\_results.tsv}.

\textbf{Effect.}  Presentation only; the release files were correct throughout.  The
tables of the 2026-09-03 build are the first to pass the automated check.

\subsection{Two central-value weight rules were in use; the fake data skipped events the prediction counted}

The framework's \texttt{safe\_weight()} sets a non-finite, negative or $>30$ central-value
weight to unity. The count-level macros (cut-and-count, cut-flow, sidebands, transfer
factors) instead zeroed invalid weights and did not clamp, and the per-run fake-data throws
\emph{skipped} events with a non-finite weight. About $1.15\%$ of overlay events carry a
non-finite tune weight, concentrated in the non-signal $\nu_\mu$CC classes, so the
pseudo-data sat below the framework's own prediction by $0.9\%$ in the signal region
($13$--$14$ events) and by $2.8$--$3.4\%$ in the CC$0\pi$ control region ($341$ FHC,
$422$ RHC POT-scaled events); the sideband yields of the note differed from the
systematics tables by the same amounts. Found through the third external review
(2026-09-03).

\textbf{Correction.} One rule, the safe-weight definition in the Monte Carlo section of the note, is now applied
everywhere: the selection's cut-flow histograms (reprocessed), every count-level macro,
and the fake-data throws, which were re-thrown with the same seeds and the events kept at
unit weight; all fifty-one universe files were rebuilt from the re-thrown fake data.
\textbf{Effect.} Pseudo-data yields rise by $0.9\%$ (signal region) to $3.4\%$ (CC$0\pi$).
Because the re-throw is a new Poisson realisation, the individual Wiener-SVD integrals
moved by up to $20\%$ (RHC $p_\mu$ $0.573\to0.690$) within their statistical
uncertainties, the inclusive closure mean moved from $1.016$ to $1.018$ (all $p>0.9$), and
the prediction-total uncertainty that the cut-and-count totals inherit from the $p_\mu$
extraction moved from $44.4\%$ to $37.7\%$ (RHC) and from $37.0\%$ to $53.5\%$
(proton-tagged RHC); the cut-and-count central values, which use the central-value
prediction rather than the throw, are unchanged. The D'Agostini cross-check was re-run on
the new release: its integral is $20$--$29\%$ above Wiener-SVD and flat across observables
to $\le1.0\%$ (previously $27$--$33\%$ and $\le2.1\%$). The real-data path was never affected.

\subsection{The total cross section borrowed its systematic from the $p_\mu$ extraction}

The headline total was a cut-and-count evaluation whose systematic fraction was taken from
the prediction total of the released $p_\mu$ extraction of the same configuration. That
proxy carried the regularisation amplification of the flux and detector terms, and it
depended on the fake-data throw through the unfolding operator: on the re-throw of
2026-09-03 alone it moved from $44.4\%$ to $37.7\%$ (RHC) and from $37.0\%$ to $53.5\%$
(proton-tagged RHC) with no change to the central value.

\textbf{Correction.} The total is now extracted with the framework as a one-bin
distribution: one true bin containing every signal event, defined on $\cos\theta_\mu$
over its full range so that no overflow exists, one reconstructed bin, Wiener filter off.
The response is the efficiency alone, nothing is regularised, $A_C\equiv1$ exactly, and
the covariance is propagated universe by universe like every other extraction. Six new
universe files (inclusive and proton-tagged $\times$ FHC/RHC/combined; SLURM 3228639) and
\texttt{data\_release/total\_xsec.tsv}. The cut-and-count macro is kept as the
closed-form check and reproduces the one-bin extraction of the central-value prediction to
four digits.
\textbf{Effect.} Inclusive totals $1.03\pm0.40$, $0.93\pm0.35$, $0.97\pm0.37$
(previously $\pm0.42$, $\pm0.35$, $\pm0.34$ with the borrowed systematic); proton-tagged
$0.62\pm0.26$, $0.53\pm0.25$, $0.57\pm0.24$ (previously $0.63\pm0.32$, $0.55\pm0.30$,
$0.59\pm0.28$). The central values now fluctuate with the throw (closure $0.98$--$1.00$
against the realised truth) instead of being the injected tune by construction. The
generator comparison moves from ``$0.5$--$0.6\sigma$ above GENIE, within $0.7\sigma$ of
the rest'' to ``$0.5$--$0.8\sigma$ above GENIE, within $0.8\sigma$ of the rest''.

\subsection{Corrections from the independent review of 2026-09-04}

Two reviewers with no prior contact with the analysis read Draft~1.4 in full, one checking
every number against the release files and one as a referee. What they found, and what
changed:

\textbf{Release index.} \texttt{data\_release/index\_extractions.tsv} and \texttt{README.md}
were the 2026-08-31 files and disagreed with every other release artefact; regenerated
from the current results (\texttt{report/tools/index\_extractions.py}), with the
differential $W_{\pi p}$ shape flagged withdrawn.

\textbf{Proton-tagged cut-flow.} The supplement's proton-tagged cut-flow table predated
the beam-off gate and PPFX-weight corrections, and the note's proton-tagged prose quoted
it ($47$--$48\%$ purity, $33\to48\%$). The fourteen overlays and the dirt overlay were
reprocessed with the proton-tagged selection under the unified weight rule
(\texttt{processed/cf\_1p}); the macro's final stage, which had counted events
unweighted ($2.6\%$ high), now carries the same weight as the other stages, and its
final row reproduces the cut-and-count check exactly ($376.4$ signal, $723.0$ predicted,
FHC). Purity is $52\%$ (FHC), $50\%$ (RHC), $51\%$ (combined); the tag rejects $67\%$ of
the neutrino background for a $32\%$ signal loss and reduces the beam-off by
$3.3$--$3.6$ at the final cut, not six.

\textbf{Flux term.} The note attributed the $24$--$33\%$ unfolded flux term to amplification
by the unfolding operator. The one-bin total, which has no regularisation, carries
$28$--$30\%$: a flux universe scales signal and background together, so the term is
$\delta\Phi/\Phi\times(1+B/S)$, $17.7\%\times1.71=30.2\%$ against the one-bin $30.25\%$
(FHC). The section is rewritten; the reconstructed-level table of the earlier pass is
dropped.

\textbf{Control-region procedure.} The target class of the $\pi^0$ region as tabulated
($\nu_\mu$CC with a $\pi^0$) is $46\%$ of the region, below the $50\%$ purity the
procedure requires for a fitted correction; stated, with the CC$+$NC~$\pi^0$ definition
to be ratified. The pass/fail criterion is now a global normalisation $\chi^2$ over the
eight (region, mode) totals; the sixteen independent $2\sigma$ and $p>0.05$ tests of the
earlier draft (which would have flagged at least one region by chance more often than
not) are diagnostics, reported with pulls against the statistical$+$detector covariance
and a normalisation-marginalised shape test. The response to a localised reconstruction
failure, the joint FHC/RHC treatment and the leverage of the held-fixed signal
contamination are specified. A control-region skim that never writes signal-region
events is added to the prerequisites; it does not yet exist.

\textbf{Ratified 2026-09-06.} The analysers ratified the $\pi^0$ region's target class as
CC$+$NC~$\pi^0$ (purity $50.1\%$ FHC, $51.0\%$ RHC of the neutrino MC; $T=0.019$--$0.020$;
transfer uncertainty $17$--$19\%$, floor-limited); the class was added to the transfer,
transfer-uncertainty and detector-variation macros and to the note's tables, and the
purity denominator of the procedure is now defined as the neutrino-MC prediction.

\textbf{Detector term statistics.} The Monte-Carlo-statistical floor of the
detector-variation term was measured on split samples
(\texttt{macros/sb\_detvar\_split.C}): the CC$0\pi$ transfer term ($27\%$ FHC, $31\%$
RHC) is genuine (floor $9\%$, $6\%$); the $\pi^0$ transfer term ($15\%$, $11\%$) is at or
below its floor ($13\%$, $17\%$) and is an upper limit; the region-yield terms have floors
below $2\%$. The procedure now uses the larger of the term and its floor.

\textbf{Ensemble.} The 32-throw ensemble was thrown by a macro that still skipped
non-finite-weight events, which accounts for $\approx1.4\%$ of its $2.5\%$ negative offset.
The macro now applies the unified rule and the ensemble was re-run for $p_\mu$,
$\cos\theta_\mu$ and the one-bin total (SLURM 3277674, 3277691, 3277692; statistical
covariances re-harvested with \texttt{slurm/ens\_statcov.sh}, pulls with
\texttt{report/tools/ensemble\_stat\_pulls.py}). Statistics-only pull widths $1.08$,
$0.97$, $1.02$ (previously $1.04$, $0.96$); ensemble means below the reference by $0.5\%$,
$2.7\%$ and $1.0\%$ at $0.4$, $1.9$ and $1.3\sigma$ (previously $2.5\%$ at $2.1\sigma$
for $p_\mu$). The bias statement in the note is reworded accordingly: a small common
negative offset, not significant, largest in the most regularised observable, and present
at the $1\%$ level in the unregularised total.

\textbf{Also.} The $A_C$ claims made about proton-tagged closure ratios (which are ratios to
the $A_C$-smeared truth) are removed; the RHC-below-FHC comparison of $A_C$-suppressed
integrals is replaced by the one-bin ratio; the reproducibility claim is scoped to the
2026-09-02 pass; the D'Agostini cross-check is stated to be prior-dominated on
central-value input; the absence of an alternative-generator fake-data test on the
released chain, of an MCS momentum systematic and of gate-ratio and Run-2 stand-in
beam-off terms is stated in the Limitations; the signal definition (one charged pion of
any momentum, above threshold), the unfolder provenance, the regularisation matrix and
the proton-tag values ($\mathrm{LLR}<0.05$, $p_p>0.3$~GeV/$c$) are stated precisely; the
flux mean energies ($0.48$~GeV) and the fraction above $0.3$~GeV ($40\%$) are given, which
is the origin of the factor $\sim\!3$ against the BNB total. Some twenty further inline
numbers that had survived from older passes were corrected.

\subsection{Documentation-only corrections from the 2026-09-02 review}

None of the following moved a released number; they are recorded because a reader of
Draft~1.1 would otherwise take the printed formulas as the implementation.
(i)~The differential cross-section equation divided by the efficiency in addition to
unfolding with the smearceptance matrix $S=\varepsilon M$; the implementation applies
the correction once, through $S$, and the equation now says so.
(ii)~The response-matrix figure was described as the column-normalised migration
matrix; it is $S$, with columns summing to the efficiency.
(iii)~The muon-candidate cut was described as ``longest qualifying track''; the code
orders the qualifying set by the muon-BDT score.
(iv)~Coherent pion production was described as a background; the signal definition is
topological and includes it.
(v)~The single-exposure normalisation example ($7.63\times10^{20}$~POT) and a
pre-rebuild systematic-breakdown table were replaced by the final three-configuration
values.  (vi)~Flux composition was labelled as integrated above $0.1$~GeV; the
values are for $60$~MeV.  (vii)~TKI angle units in the binning appendix read rad; they
are degrees.  (viii)~A $\theta_\pi$ result was listed as final; none is extracted.
(ix)~The beamline-geometry flux term, quoted at $1.6$--$3.3\%$, is not in the released
systematics configuration and is now stated as absent.

\section{Development-era corrections}

These predate the framework-based analysis and are recorded for completeness.  They were
previously listed in a ``Known Issues'' table in the analysis note; that table had become
thirteen resolved rows against one open one, so it was retired and its contents moved
here, where a resolved correction belongs.

\begin{center}
\small
\begin{tabular}{p{0.36\textwidth}p{0.56\textwidth}}
\toprule
Item & Resolution \\
\midrule
Unfolding covariance was statistics-only
  & Replaced by the framework \texttt{PredTotal} plus the data Poisson term, removing a
    ${\sim}26\%$ bias in the quoted uncertainty. \\
$\theta_{\mu\pi}$ binning mismatch
  & A custom 10-bin scheme was aligned to the framework's 9-bin one; the two pipelines had
    been binning the same observable differently. \\
Negative tail bins
  & An artefact of the EXT treatment; all bins positive once the subtraction was
    corrected. \\
EXT subtraction on fake data
  & The custom pipeline cancelled it via \texttt{fake\_data\_ext\_cancel}, a flag that had
    to be remembered at unblinding. Superseded: the framework selects the beam-off
    treatment from the presence of Monte Carlo truth in the beam-on file (added to
    pseudo-data, subtracted from real data) and refuses a real file unless
    \texttt{XSEC\_UNBLIND=1}, so no manual switch remains. \\
Multisim propagation in the custom pipeline
  & Never completed, and now moot: the released results come from the framework, which
    propagates the multisims (\texttt{cov\_xsec\_multi} is in the release).  The row
    survived in the note as the single ``open'' item long after it had ceased to describe
    anything. \\
CRT veto branches
  & \texttt{crthitpe} and \texttt{crtveto} are all-zero in the NuMI ntuples; replaced by
    \texttt{bdt\_cosmic}, disabled at its default. \\
\bottomrule
\end{tabular}
\end{center}

\section{Open items}

Two things remain outstanding, and neither is a correction awaiting application: they are
limits of the present exposure and of the blinding policy.

\paragraph{The differential $W_{\pi p}$ shape.}
Withdrawn, and not recoverable by any systematic treatment.  At this exposure no binning
satisfies both the migration and the statistics requirement simultaneously: a two-bin
split near $1.15$~GeV/$c^2$ meets the response criterion but leaves about $137$ events in
the first bin, which the Wiener filter then suppresses, while wider bins restore the
statistics and fail the response.  The proton-tagged sample supports an integrated cross
section.

\paragraph{External reproduction of the $A_C$-smeared comparison --- resolved.}
This was listed here as open in the previous revision.  It is resolved.  The transformation
\[
  \left(d\sigma/dx\right)^{\rm smeared}_i = \frac{1}{w_i}\sum_j (A_C)_{ij}\,p_j ,
\]
applied to a prediction supplied as a per-bin cross section, reproduces every released
generator curve to a worst per-bin deviation of $1.2\times10^{-10}$.  It was the failure of
this very test that exposed the double-$A_C$ bug, and its success is now the regression
check against recurrence.

What remains is narrower and is not a numerical question: no packaged reference
implementation ships with the release, so a user must write the three lines above
themselves.  They can verify their implementation against any model column of
\texttt{curves\_*.tsv}.

\paragraph{Everything that requires unblinding.}
No beam-on signal region has been opened, so the control-region checks in data, the
background constraints, and any statement that a result agrees or disagrees with a
generator remain untested against real data.  Every number in the analysis note is a
property of the injected sample and the extraction machinery.

\paragraph{A caveat that is not an open item but reads like one.}
The closure $\chi^2$ values have $p$-values clustered near unity (median $0.98$ over the
fifty-one extractions) because the measurement and the truth it is compared against share
their systematics, so those terms cancel in the difference while still inflating the
covariance used to judge it.  Sizing the uncertainty band properly would need a
restricted covariance or an ensemble of throws.  This is a limitation of what the closure
test demonstrates, not an error in the values.

\section{Superseded values that may still be encountered}

\begin{center}
\begin{tabular}{lll}
\toprule
Quantity & superseded value & status \\
\midrule
Mean inclusive closure ratio        & $1.106$, $1.001$, $1.016$ & now $1.018$ (truth-side correction; re-thrown fake data) \\
RHC $W_{\pi p}$ cross section       & $-0.194$         & sign was a POT-convention artefact \\
Combined detector systematic        & $42.3\%$         & now $24.6\%$ \\
Combined result / expected          & $0.910$          & now $1.027$ \\
FHC $\theta_\mu$ $A_C$ conditioning & $0.015$          & now $0.51$ \\
Mean $\delta p_T$                   & $0.742$~GeV/$c$  & now $0.264$~GeV/$c$ (beam frame) \\
Proton-tagged \texttt{detVar}       & \texttt{n/a}, then $10.5$--$64.4\%$ & now $10.5$--$108.5\%$ (median $33.0\%$, 2026-09-04 release) \\
Proton-tagged $\sigma$, all obs.    & 2026-08-18 values & superseded 2026-08-31 \\
\bottomrule
\end{tabular}
\end{center}

\end{document}
